% !TEX root = ../main.tex
\chapter{The Autograd Engine}
\label{ch:autograd_engine}

The previous chapter explained reverse-mode AD as an algorithm.
This chapter explains \Lucid's implementation of it. The autograd
engine lives in \code{lucid/\_C/autograd/} (\chref{ch:layered_dag},
Layer 4) and is the component that constructs the backward graph
during forward execution, walks it during \code{backward()}, and
accumulates cotangents into parameter \code{.grad} fields.

The engine is conventional in shape --- reverse-mode AD engines
have converged on a similar structure across frameworks --- but
the details are where the implementation lives. We treat the
\code{Node} class, the engine's traversal algorithm, the
saved-tensor mechanism, the version-counter check that catches
in-place modification errors, the hook system, and the anomaly
detector.

\secref{sec:autograd_engine:overview} sketches the architecture.
\secref{sec:autograd_engine:node} treats the \code{Node} class.
\secref{sec:autograd_engine:construction} treats forward-time
graph construction.
\secref{sec:autograd_engine:traversal} treats backward traversal.
\secref{sec:autograd_engine:accumulate} treats
\code{AccumulateGrad} on leaves.
\secref{sec:autograd_engine:saved} treats saved tensors.
\secref{sec:autograd_engine:version} treats the version counter.
\secref{sec:autograd_engine:hooks} treats hooks.
\secref{sec:autograd_engine:anomaly} treats the anomaly detector.

\section{Architectural Overview}
\label{sec:autograd_engine:overview}

The engine has three principal components, each implemented as a
C++ class under \code{lucid/\_C/autograd/}:

\begin{itemize}
  \item \code{Node}: one node in the backward graph. Concrete
    subclasses (one per op category) implement the VJP.
  \item \code{Engine}: the traversal driver. Walks the backward
    graph in reverse topological order, calling each node's
    \code{apply} with the accumulated cotangent.
  \item \code{AccumulateGrad}: a special leaf node that
    accumulates the final cotangent into a tensor's \code{.grad}
    field.
\end{itemize}

The supporting machinery includes the \code{AutogradMeta} on each
tensor (\chref{ch:tensor_model}~Section~\ref{sec:tensor_model:autograd_meta}),
the saved-tensor list per node, the hook registries, and the
anomaly detector.

\figref{fig:autograd_engine:overview} sketches the relationship.

\begin{figure}[t]
\centering
\begin{adjustbox}{max width=\textwidth, center}
\begin{tikzpicture}[
  box/.style={draw, rounded corners=2pt, font=\small\sffamily,
              minimum width=3.5cm, minimum height=1.1cm,
              align=center, inner sep=4pt},
  tensor/.style={box, fill=lucid.note.bg},
  node/.style={box, fill=lucid.info.bg},
  leaf/.style={box, fill=lucid.ok.bg},
  engine/.style={box, fill=lucid.code.bg, minimum width=4cm,
                 minimum height=1.4cm},
  fwd/.style={->, >=stealth, thick, color=lucid.accent.dark,
              shorten >=2pt, shorten <=2pt},
  bwd/.style={->, >=stealth, thick, color=lucid.warn.fg,
              dashed, shorten >=2pt, shorten <=2pt},
]
  % Top: tensors
  \node[tensor] (param) at (-5, 3) {parameter\\(leaf, \texttt{requires\_grad})};
  \node[tensor] (act1) at (0, 3) {intermediate};
  \node[tensor] (loss) at (5, 3) {loss\\(scalar)};

  % Middle: backward nodes
  \node[leaf] (accumg) at (-5, 0.5) {AccumulateGrad};
  \node[node] (n_op1) at (0, 0.5) {MulBackward};
  \node[node] (n_op2) at (5, 0.5) {SumBackward};

  % Bottom: engine
  \node[engine] (eng) at (0, -1.8) {Engine::backward()\\(topological walk)};

  \draw[fwd] (param) -- node[font=\footnotesize, right] {grad\_fn} (accumg);
  \draw[fwd] (act1) -- node[font=\footnotesize, right] {grad\_fn} (n_op1);
  \draw[fwd] (loss) -- node[font=\footnotesize, right] {grad\_fn} (n_op2);

  \draw[bwd] (n_op2) -- (n_op1);
  \draw[bwd] (n_op1) -- (accumg);

  \draw[fwd, color=lucid.muted] (eng) -- (n_op2);
  \draw[fwd, color=lucid.muted] (eng) -- (n_op1);
  \draw[fwd, color=lucid.muted] (eng) -- (accumg);
\end{tikzpicture}
\end{adjustbox}
\caption[Autograd engine architecture.]{The autograd engine. Each
tensor that requires grad carries a \code{grad\_fn} pointer to
its backward \code{Node}. Nodes form a backward DAG (dashed
arrows) that mirrors the forward DAG. \code{Engine::backward()}
walks the backward DAG in reverse topological order, applying
each node's VJP. Leaves use \code{AccumulateGrad} to deposit
their cotangent into the tensor's \code{.grad} field.}
\label{fig:autograd_engine:overview}
\end{figure}

\section{The Node Class}
\label{sec:autograd_engine:node}

\code{Node} is the base class for all backward operations.
Defined in \code{lucid/\_C/autograd/Node.h}:

\begin{lstlisting}[style=lucidcpp,language=C++]
class Node {
public:
  virtual ~Node() = default;

  // Apply this node's VJP: take output cotangents, return
  // input cotangents.
  virtual std::vector<TensorImpl> apply(
      const std::vector<TensorImpl>& output_grads) = 0;

  // The list of upstream nodes that this node will propagate
  // gradient to (one per input that requires grad).
  std::vector<std::shared_ptr<Node>> next_nodes;

  // For each next_node, which output position of that node
  // does this node's input correspond to.
  std::vector<uint32_t> next_output_indices;

  // The forward inputs / outputs that were saved for this
  // backward.
  std::vector<TensorImpl> saved_tensors;

  // Hooks
  std::vector<Hook> pre_hooks;
  std::vector<Hook> post_hooks;
};
\end{lstlisting}

\subsection{Per-op subclasses}
\label{ssec:autograd_engine:per_op_subclass}

Each differentiable op has a corresponding \code{Node} subclass
that implements \code{apply}. The pattern is uniform:

\begin{lstlisting}[style=lucidcpp,language=C++]
class MulBackward : public AutogradNode<MulBackward, 2> {
public:
  // saved_tensors = {a, b}; output_grads = {grad_y}
  std::vector<TensorImpl> apply(
      const std::vector<TensorImpl>& output_grads) override {
    const auto& grad_y = output_grads[0];
    const auto& a = saved_tensors[0];
    const auto& b = saved_tensors[1];
    return {
      Dispatcher::instance().binary(OpKind::Mul, b, grad_y),
      Dispatcher::instance().binary(OpKind::Mul, a, grad_y),
    };
  }
};
\end{lstlisting}

\code{AutogradNode<Derived, NInputs>} is a CRTP base that
provides boilerplate (\code{next\_nodes} sizing, default
\code{saved\_tensors} handling). Per-op subclasses provide only
the \code{apply} body.

\subsection{Why per-op classes, not function pointers}
\label{ssec:autograd_engine:why_classes}

A function-pointer approach would have been simpler in line
count but loses two properties:

\begin{itemize}
  \item \textbf{Polymorphic state.} Some backwards need state
    beyond the saved tensors (e.g., \code{ConvBackward} needs the
    stride and padding from the forward call). Per-op subclasses
    hold this state cleanly; function pointers would have needed
    a void-pointer context.
  \item \textbf{Debugging.} A \code{Node} subclass has a name that
    appears in stack traces and the anomaly detector's output.
    Function pointers would have lost this.
\end{itemize}

The boilerplate cost is amortised by the CRTP base.

\section{Forward-Time Construction}
\label{sec:autograd_engine:construction}

During the forward pass, when a differentiable op runs, the
dispatcher consults the op's schema and, if backward is needed,
constructs a backward node and attaches it to the output tensor's
\code{AutogradMeta}.

\subsection{The construction sequence}
\label{ssec:autograd_engine:construction_seq}

For \code{c = a * b} where \code{a.requires\_grad == True}:

\begin{enumerate}
  \item Dispatcher routes to backend, computes \code{c} value.
  \item Dispatcher checks: does any input require grad? If yes,
    proceed.
  \item Construct \code{MulBackward} node.
  \item Save \code{a} and \code{b} into the node's
    \code{saved\_tensors}.
  \item Wire \code{next\_nodes}: \code{next\_nodes[0] =
    a.grad\_fn}, \code{next\_nodes[1] = b.grad\_fn} (or
    \code{AccumulateGrad} for leaves).
  \item Set \code{c.grad\_fn = MulBackward instance}.
  \item Set \code{c.requires\_grad = True} (inherits from
    inputs).
  \item Return \code{c}.
\end{enumerate}

The whole sequence takes a few microseconds on top of the actual
multiplication. For multiplications on small tensors the
overhead is noticeable; for any non-trivial op the kernel
dominates.

\subsection{Disabling under no\_grad}
\label{ssec:autograd_engine:no_grad}

The construction sequence above is skipped when the global
\code{GradMode} flag is off (\code{lucid.no\_grad()} context).
\code{c.grad\_fn} remains null, \code{c.requires\_grad} remains
false, and no autograd state is constructed.

This is the standard pattern for evaluation: wrap the inference
forward in \code{with lucid.no\_grad():} to avoid building a
backward graph that will never be traversed.

\subsection{Inference mode}
\label{ssec:autograd_engine:inference_mode}

\code{lucid.inference\_mode()} is a stricter version of
\code{no\_grad}: it additionally disables the version counter
and a few other autograd hooks, giving slightly better
performance at the cost of forbidding any later differentiation
through tensors created inside the context.

The cost difference is roughly 5--10\,\% on small-tensor
workloads; \code{no\_grad} is the safer default.

\section{Backward Traversal}
\label{sec:autograd_engine:traversal}

When the user calls \code{loss.backward()}, the engine walks the
backward graph in reverse topological order, calling each node's
\code{apply}.

\subsection{The traversal algorithm}
\label{ssec:autograd_engine:algorithm}

The algorithm is standard reverse-mode AD:

\begin{lstlisting}[style=lucidcpp,language=C++]
void Engine::backward(TensorImpl& loss) {
  // Seed the queue with the loss's grad_fn
  auto root = loss.grad_fn();
  std::unordered_map<Node*, std::vector<TensorImpl>> cotangents;
  cotangents[root.get()] = {ones_like(loss)};

  // Topological sort
  auto order = topological_sort_from(root);

  // Walk in reverse order, applying each node
  for (auto& node : order) {
    auto& outputs = cotangents[node.get()];
    auto inputs = node->apply(outputs);
    // Distribute input cotangents to upstream nodes
    for (size_t i = 0; i < node->next_nodes.size(); ++i) {
      auto& upstream = node->next_nodes[i];
      uint32_t pos = node->next_output_indices[i];
      cotangents[upstream.get()].at(pos) =
          add_or_accumulate(cotangents[upstream.get()].at(pos),
                            inputs[i]);
    }
  }
}
\end{lstlisting}

The key step is the accumulation: if a node has multiple
upstream contributors (e.g., the multi-variable chain rule's
$+=$), each contribution adds rather than replaces.

\subsection{Topological sort}
\label{ssec:autograd_engine:topo_sort}

The topological sort uses Kahn's algorithm: identify nodes with
no outgoing edges (the loss's \code{grad\_fn} initially), emit
them, decrement the in-degree of their predecessors, and repeat.

The sort is done once per \code{backward()} call. For a graph
with $N$ nodes it costs $O(N)$ time. For a typical training step
with a few thousand nodes, the sort itself is negligible
relative to the per-node \code{apply} cost.

\subsection{Multiple roots}
\label{ssec:autograd_engine:multi_root}

\code{Engine::backward()} accepts a list of root tensors and a
list of seed gradients (one per root). For most uses there is one
root (the loss) and one seed (\code{ones\_like(loss)}). The
multi-root form is used for vector-valued objectives like
multi-task loss, where the backward is the weighted sum of per-
task backwards.

\section{AccumulateGrad on Leaves}
\label{sec:autograd_engine:accumulate}

A leaf tensor (parameter created by the user with
\code{requires\_grad=True}) does not have a forward op that
produced it, so it has no per-op backward node. Instead, the
engine attaches an \code{AccumulateGrad} node to it.

\subsection{The AccumulateGrad node}
\label{ssec:autograd_engine:accumg_node}

\code{AccumulateGrad} is the simplest \code{Node}: its
\code{apply} adds the incoming cotangent to the leaf tensor's
\code{.grad} field. There are no upstream nodes; the cotangent
stops here.

\begin{lstlisting}[style=lucidcpp,language=C++]
class AccumulateGrad : public Node {
public:
  TensorImpl variable;  // the leaf tensor

  std::vector<TensorImpl> apply(
      const std::vector<TensorImpl>& output_grads) override {
    const auto& g = output_grads[0];
    if (!variable.autograd_meta()->grad) {
      variable.autograd_meta()->grad = g.clone();
    } else {
      Dispatcher::instance().binary(
          OpKind::Add,
          *variable.autograd_meta()->grad,
          g);
    }
    return {};  // no upstream
  }
};
\end{lstlisting}

\subsection{The optimizer.zero\_grad pattern}
\label{ssec:autograd_engine:zero_grad}

Because \code{AccumulateGrad} adds rather than replaces, gradients
accumulate across \code{backward()} calls until explicitly
cleared. The standard pattern at the start of each training step:

\begin{lstlisting}[style=lucid,language=LucidPython]
optimizer.zero_grad()
loss = model(x)
loss.backward()
optimizer.step()
\end{lstlisting}

\code{optimizer.zero\_grad()} iterates over the optimiser's
parameter groups and sets each \code{.grad} to None (which causes
the next \code{AccumulateGrad} to clone rather than add). The
default behaviour matches \refframework's; \code{set\_to\_none=False}
zeros in place if memory pressure forbids re-allocation.

\section{Saved Tensors}
\label{sec:autograd_engine:saved}

The backward function for an op needs forward quantities to
compute its VJP. The framework saves these quantities at forward
time and holds them until backward consumes them.

\subsection{What gets saved}
\label{ssec:autograd_engine:what_saved}

The op's schema declares which forward tensors are needed:

\begin{itemize}
  \item Some backwards need the inputs (e.g., \code{mul} needs
    both inputs).
  \item Some need the output (e.g., \code{sigmoid} backward uses
    $y(1-y)$, where $y$ is the output).
  \item Some need both.
  \item Some need neither (e.g., \code{add} backward is just the
    identity per input).
\end{itemize}

The schema's \code{BackwardSpec::saved} list names the saved
quantities. The engine snapshots them into the node's
\code{saved\_tensors} vector at forward time.

\subsection{Memory cost}
\label{ssec:autograd_engine:saved_memory}

Saved tensors are the dominant memory cost of training. For a
deep model with $L$ layers and per-layer activation size $S$, the
total saved memory is $L \cdot S$. This grows linearly with depth
and is the reason gradient checkpointing
(\chref{ch:gradient_checkpointing}) exists.

\subsection{Why save outputs when possible}
\label{ssec:autograd_engine:save_outputs}

For ops like \code{sigmoid} and \code{tanh}, the backward can be
computed from either the input or the output (since the output is
a function of the input). Saving the output is preferable because
the output is computed anyway and saving it costs no additional
forward computation; saving the input would require keeping the
input alive after the forward, which the runtime might otherwise
have freed.

The schema's \code{BackwardSpec::saved} for these ops explicitly
names the output, and the engine acts accordingly.

\section{The Version Counter}
\label{sec:autograd_engine:version}

In-place modification of a tensor that was saved for backward
would silently corrupt the gradient: the backward function would
read the modified value instead of the original. The version
counter catches this class of bug.

\subsection{Mechanism}
\label{ssec:autograd_engine:version_mech}

Each \code{TensorImpl} carries a 64-bit \code{version\_counter}
that increments on every in-place op (\code{add\_}, \code{mul\_},
\code{copy\_}, etc.). When a tensor is saved for backward, the
engine captures the current counter value. At backward time, the
engine compares the captured value to the current value. A
mismatch raises:

\begin{verbatim}
LucidError: one of the variables needed for gradient computation
has been modified by an inplace operation. The variable was a
saved tensor of MulBackward, which expected version 0 but found
version 1.
\end{verbatim}

\subsection{When the check fires}
\label{ssec:autograd_engine:version_fires}

The check fires for code like:

\begin{lstlisting}[style=lucid,language=LucidPython]
x = lucid.tensor([1.0, 2.0], requires_grad=True)
y = x * 2          # MulBackward saves x
x.add_(1.0)        # in-place: x.version_counter goes 0 -> 1
y.sum().backward() # raises: version mismatch on saved x
\end{lstlisting}

The user fix is either to avoid the in-place op or to clone
first: \code{x = x.clone(); x.add\_(1.0); ...}.

\subsection{When the check doesn't fire}
\label{ssec:autograd_engine:version_misses}

The version counter is per-\code{TensorImpl}, not per-byte. An
in-place op on a view of the saved tensor modifies the underlying
storage but not the saved tensor's version counter; the check
misses. The user should avoid in-place ops on views of saved
tensors.

This is one of the few correctness gaps in the autograd engine.
The mitigation is convention (avoid in-place on saved tensors)
plus the anomaly detector (which catches NaN gradients that often
indicate corruption).

\section{Hooks}
\label{sec:autograd_engine:hooks}

The user can attach \emph{hooks} to tensors and modules to
intercept the autograd flow without modifying the op
implementation.

\subsection{Tensor hooks}
\label{ssec:autograd_engine:tensor_hooks}

\code{tensor.register\_hook(fn)} attaches a function to a tensor's
\code{grad\_fn}. The function is called with the cotangent before
it propagates upstream, and may return a modified cotangent.

Use cases:
\begin{itemize}
  \item Gradient clipping at a specific tensor (rare; usually
    done at the parameter level).
  \item Logging or visualisation (sniff the cotangent without
    changing it).
  \item Gradient masking (zero out cotangents for certain
    components).
\end{itemize}

\subsection{Module hooks}
\label{ssec:autograd_engine:module_hooks}

\code{module.register\_forward\_hook(fn)} and
\code{register\_backward\_hook(fn)} attach to a \code{Module}
instance and fire on every forward and backward through the
module. The forward hook sees the module inputs and outputs; the
backward hook sees the gradient inputs and outputs.

The backward hook is the typical mechanism for debugging
gradient flow: print the magnitude of the gradient at each
layer to find vanishing or exploding gradients.

The chapter for module hooks proper is
\chref{ch:module_system}.

\section{The Anomaly Detector}
\label{sec:autograd_engine:anomaly}

\code{lucid.autograd.set\_detect\_anomaly(True)} enables a
diagnostic mode that helps locate the source of NaN or Inf in
gradients.

\subsection{What it does}
\label{ssec:autograd_engine:anomaly_does}

When enabled:
\begin{itemize}
  \item The forward pass records the Python call stack of every
    op for later use in error messages.
  \item The backward pass checks each output cotangent for NaN
    or Inf. The first occurrence raises an error with the saved
    forward call stack, identifying the originating op.
\end{itemize}

Without the anomaly detector, a NaN in a gradient propagates
silently and the user sees only the eventual symptom (NaN
parameters after the optimiser step, training divergence). The
detector trades performance (5--10x slowdown of the backward
pass) for diagnostic clarity.

\subsection{When to use}
\label{ssec:autograd_engine:anomaly_when}

Recommended whenever training mysteriously produces NaN
gradients. Not recommended for normal training; the slowdown is
substantial.

\section{Summary and Forward References}
\label{sec:autograd_engine:summary}

\Lucid's autograd engine implements reverse-mode AD through a
\code{Node}-based backward graph constructed at forward time and
traversed by \code{Engine::backward()} in reverse topological
order. Per-op \code{Node} subclasses implement the VJP;
\code{AccumulateGrad} handles leaf accumulation; the version
counter catches in-place corruption; hooks expose the autograd
flow to user code; the anomaly detector helps diagnose NaN
gradients.

The next chapter, \chref{ch:custom_functions}, treats the
user-facing \code{Function} class that lets the user define
custom ops with custom backwards.
\chref{ch:gradient_checkpointing} treats the memory-saving
recomputation pattern. \chref{ch:higher_order} treats functional
transforms.

\begin{lucidnote}
For the user: \code{loss.backward()} is the only entry point.
Behind it, the engine walks the saved DAG and accumulates
gradients into \code{.grad}. The version counter, the hook
system, and the anomaly detector are the three knobs the user can
use to inspect and intervene in the process.
\end{lucidnote}
